\chapter{Conclusiones y trabajo futuro}
\label{chap:conclusiones}

El trabajo abordó el desarrollo de un sistema de visión por computadora para obtener información cuantitativa del movimiento de microorganismos acuáticos a partir de videos adquiridos mediante microscopía óptica en una línea experimental del GFA. Este capítulo sintetiza los aportes de la implementación y de su evaluación, y presenta las líneas que permitirían ampliar su aplicación experimental. La discusión de los resultados y de sus limitaciones se desarrolló en el capítulo~\ref{chap:evaluacion-resultados}.

\section{Conclusiones}

Se integraron en un mismo procedimiento la detección mediante un modelo preentrenado adaptado al material disponible, la asociación temporal de las detecciones, la selección de trayectorias y el cálculo de magnitudes de movimiento. Sobre las posiciones registradas se incorporaron visualizaciones, herramientas de inspección individual y una clasificación de patrones basada en reglas. Esta integración permite reutilizar las trayectorias en distintos análisis y conservar los parámetros principales del procesamiento, sin volver a ejecutar la inferencia para cada consulta.

El aporte del sistema se sitúa en la adaptación de ese procedimiento a los videos y necesidades de esta línea experimental. La representación de cada trayectoria mediante cuadros y posiciones hace explícita la información utilizada para calcular distancias, rapidez y continuidad temporal. La separación de responsabilidades entre seguimiento, análisis e inspección facilita revisar tanto los resultados como las decisiones de cálculo que los producen. Los productos generados ofrecen una base organizada para la revisión posterior por parte de los investigadores.

Las pruebas controladas aportaron evidencia de consistencia numérica y de integración en los casos cubiertos. En particular, se verificó que la rapidez utiliza el tiempo real entre observaciones cuando existen cuadros ausentes. Este resultado es relevante porque tratar como consecutivas posiciones separadas por un gap modificaría las magnitudes temporales. La conservación de los índices originales permite además identificar los intervalos cuya información está incompleta y distinguirlos del recorrido efectivamente observado.

La evaluación experimental vinculó una trayectoria automática con una referencia manual parcial y permitió reconstruir sus productos de análisis. Después de la alineación documentada, los puntos posteriores al anclaje presentaron una separación media aproximada de $5{,}33$ px. La comparación respalda la similitud espacial del caso examinado. El contraste de rapidez mostró, a su vez, que esa similitud debe analizarse junto con las discrepancias de las magnitudes derivadas y con la resolución temporal de las anotaciones. La validación de posiciones y la de rapidez representan, por ello, aspectos complementarios de la evaluación.

Los resultados de entrenamiento respaldan la viabilidad de adaptar el detector al material conservado. Su interpretación se mantiene asociada a las dos experiencias realizadas y a las evaluaciones registradas. La comparación detallada de un individuo aporta evidencia sobre un caso real, mientras que una caracterización global del seguimiento requiere referencias adicionales que permitan examinar asociaciones, cruces y fragmentación en el conjunto del video.

La clasificación incorpora una descripción computacional inspeccionable de las trayectorias. Sus puntuaciones permiten revisar la decisión obtenida y relacionarla con las características del recorrido. El alcance comprobado corresponde a la aplicación consistente de las reglas implementadas; su adecuación para distinguir comportamientos de interés experimental requiere categorías de referencia independientes. Del mismo modo, los umbrales de inactividad y de cambios de rapidez conservan un carácter operativo.

En conjunto, el trabajo proporciona una implementación y evidencia de su funcionamiento sobre datos controlados y un caso experimental acotado. El sistema puede apoyar la organización y revisión de mediciones del movimiento, manteniendo la intervención de los especialistas para interpretar su significado físico o biológico. La comparación cuantitativa del esfuerzo respecto del procedimiento manual queda pendiente. Las magnitudes presentadas se conservan en px, px/s y px/s$^2$; la calibración espacial y la conversión a unidades físicas permanecen \textbf{A DEFINIR}.

\section{Trabajo futuro}

Las siguientes líneas se organizan según su relación con los límites identificados durante la evaluación. Su desarrollo permitiría ampliar la evidencia y precisar las condiciones de uso del sistema.

\begin{enumerate}
  \item \textbf{Calibración espacial y conversión a unidades físicas: A DEFINIR.} La definición del procedimiento, el factor de conversión y su aplicación a las mediciones quedan pendientes de resolución con la dirección del trabajo.

  \item \textbf{Ampliación de la evaluación del seguimiento.} Incorporar referencias de varios individuos y videos, con revisión de las identidades antes y después de cruces, oclusiones e interrupciones. Esta ampliación permitiría evaluar la conservación de identidad y la fragmentación, además de las discrepancias de posición y rapidez. La comparación debería conservar la correspondencia temporal y documentar los criterios de anotación y revisión manual.

  \item \textbf{Diversificación del material del detector.} Reunir nuevas imágenes originales cuando la disponibilidad de muestras lo permita, incorporando variaciones de contraste, iluminación, densidad y condiciones de adquisición. La separación de entrenamiento y evaluación debería considerar la procedencia por video o muestra para examinar el desempeño con material independiente. El efecto de nuevos datos, aumentos y configuraciones requeriría mediciones comparables y repeticiones de entrenamiento.

  \item \textbf{Evaluación de la clasificación y de los umbrales analíticos.} Construir, con participación de especialistas, un conjunto de trayectorias categorizadas mediante criterios externos al algoritmo. Esa referencia permitiría examinar los aciertos, las confusiones y los casos ambiguos, y revisar la separación entre puntuaciones y las condiciones de rechazo. También resultaría necesario estudiar la sensibilidad al muestreo, al suavizado y a los umbrales utilizados para identificar baja movilidad y cambios de régimen.

  \item \textbf{Medición del esfuerzo de análisis.} Comparar el procedimiento manual y el automático sobre los mismos videos y con productos finales equivalentes. La medición debería incluir preparación, procesamiento, corrección e inspección, distinguiendo el tiempo de ejecución del tiempo de intervención del operador. Este diseño permitiría cuantificar la utilidad práctica que motivó la automatización.
\end{enumerate}

Como mejora complementaria de la implementación, convendría completar el registro de la configuración, incluyendo los parámetros del clasificador y las huellas de las entradas, y facilitar la conservación de las posiciones obtenidas. Estos cambios permitirían reforzar la trazabilidad entre ejecuciones. Su incorporación debería acompañarse de las pruebas correspondientes y de una actualización de la guía del Apéndice~\ref{ap:guia-ejecucion}.
